% META: {"id": "1SPE-SUITES-EX-013", "chapitre": "1SPE-SUITES", "type_objet": "exercice", "capacites": ["1SPE-SUITES-C4"], "capacites_codes": ["C4"], "methodes": ["M4"], "parcours": 1, "competences": ["calculer"], "duree_min": 10, "mode_creation": "ex_nihilo", "sources_inspiration": [], "parametres_sympy": {}, "fichier_tex": "chapitres/1SPE-SUITES/exercices/1SPE-SUITES-EX-013.tex", "corrige_tex": "chapitres/1SPE-SUITES/corriges/1SPE-SUITES-CO-013.tex", "status": "generated"}
% BEGIN-VERIFY
% from sympy import *
% # Somme 1 + 2 + 3 + ... + 50
% S = sum(range(1, 51))
% assert S == 1275
% # Formule : n*(n+1)/2 avec n=50
% assert 50*51//2 == 1275
% END-VERIFY
\begin{exercice}{1SPE-SUITES-EX-013}{1}{10}
On souhaite calculer la somme $S = 1 + 2 + 3 + \cdots + 50$.
\begin{enumerate}
  \item En appliquant la formule $1 + 2 + \cdots + n = \dfrac{n(n+1)}{2}$, calculer $S$.
  \item Vérifier que $S = 1275$ en remarquant que les paires $(1 + 50)$, $(2 + 49)$, $\ldots$ valent chacune $51$, et en comptant le nombre de telles paires.
\end{enumerate}
\end{exercice}
